\section*{Abstract}

\todo{Write after Sections 2--5 stabilize. Draft claim structure (each claim
below must resolve to a citation or a generated table before this TODO is
removed):}

\begin{itemize}
\item \policyengineus{} is an open-source microsimulation model of the
United States federal, state, and local tax and benefit system, implemented
as statutory rules encoded as Python formulas over a calibrated microdata
foundation (the \populace{} stack).
\evidence{policyengine\_us/programs.yaml; docs-quarto/index.qmd}
\item Rules coverage: \todo{cite exact counts of programs, parameters, and
variables once confirmed -- see Section~2}.
\item Validation: \policyengineus{} tax liabilities agree with \taxsim{}35
within \$100 for \todo{X}\% of federal and \todo{X}\% of state tax units on
the 2022 \cps{}, per the existing validation notebook.
\evidence{docs/validation/taxsim.ipynb}
\item Data: microdata is calibrated to administrative totals (\soi{}, Census,
SSA, CMS, SNAP, TANF, \cbo{} baseline aggregates) via the \populace{}
calibration surface.
\evidence{populace/packages/populace-build/src/populace/build/us/fiscal\_target\_references.json}
\item \todo{CBO/JCT/TPC comparison claim, pending Section~5.3 evidence
mapping -- currently TODO(evidence needed) as of outline drafting.}
\end{itemize}

\vspace{1em}
\noindent\textbf{Keywords:} microsimulation, tax-benefit model, validation,
TAXSIM, administrative data
